% Source-derived chapter 02: Virasoro constraints
% Original source: virasoro-constraints-moduli-sheaves-vertex-algebras
\chapter{Virasoro constraints}
\label{virasoro-constraints-moduli-sheaves-vertex-algebras-chapter-02}
\source{virasoro-constraints-moduli-sheaves-vertex-algebras}

\begin{remark}[Source section]
\label{virasoro-constraints-moduli-sheaves-vertex-algebras-chapter-02-source}
\source{virasoro-constraints-moduli-sheaves-vertex-algebras}
\source{virasoro-constraints-moduli-sheaves-vertex-algebras}
This chapter reproduces the corresponding section of the retrieved
LaTeX source, including its definitions, statements, examples, and
proofs. It has not yet been translated into Lean.
\notready
\end{remark}

% The source section title was: Virasoro constraints
\label{sec: virasoro constraints}
\source{virasoro-constraints-moduli-sheaves-vertex-algebras}

In this section, we formulate the Virasoro constraints for moduli spaces of sheaves and pairs, denoted by $M$ and $P$, respectively. These moduli spaces satisfy the numbered assumptions in Section \ref{sec: modulisheavespairs} unless otherwise mentioned. We use the notation $M_\alpha$ when the topological type $\alpha$ of the sheaves in the moduli space is relevant. 

\subsection{Supercommutative algebras}
\label{sec: supercommutative}
\source{virasoro-constraints-moduli-sheaves-vertex-algebras}
Before we move onto geometry, we note down some useful observations about freely generated \textit{supercommutative} algebras and derivations on them. Let $D_\bullet$ be a supercommutative $\BZ$-\textit{graded unital algebra} over~$\BC$ with degree $\deg(v) = i$ for any $v\in D_i$. Supercommutativity means that multiplication satisfies 
\[v\cdot w = (-1)^{|v||w|} w\cdot v\,,\]
where 
$$|v|\in \{0,1\}\,,\quad  \text{such that} \quad |v|\equiv \deg(v)\mod 2\,.$$
The unit of $D_\bullet$ is always going to be denoted by $1$ and in general, we will omit specifying it in the notation. 

A \textit{superderivation} of degree $r$ on $D_\bullet$ is a $\BZ$-graded linear map 
$$
R:D_\bullet\longrightarrow D_{\bullet+r}
$$
satisfying the \textit{graded Leibnitz rule}
$$R(v\cdot w) = R(v)\cdot w + (-1)^{r|v|}\,v\cdot R(w)\,.$$


\begin{definition}
\source{virasoro-constraints-moduli-sheaves-vertex-algebras}
\notready\label{def: supercommutativefree}
\source{virasoro-constraints-moduli-sheaves-vertex-algebras}
Let $C_\bullet$ be a $\BZ$-graded  $\BC$-vector space. We denote by $$D_\bullet = \SSym[C_\bullet]$$ the unital supercommutative algebra freely generated by $C_\bullet$. Denote by $C^\bullet$ the graded dual of $C_\bullet$. We define the dual of $D_\bullet$ as a completion of $\SSym[C^\bullet]$ with respect to the degree. More precisely, the dual is
\[D^\bullet\coloneqq \SSym\llbracket C^\bullet\rrbracket=\prod_{i\geq 0}\SSym[C^\bullet]^i\]
where $\SSym[C^\bullet]^i$ denotes the degree $i$ part of $\SSym[C^\bullet]$ with the degree induced by the one on $C^\bullet$.\footnote{We follow here the convention that the total homology is the direct sum of the homology groups in each degree, while cohomology is the product.} 
\end{definition}
Given a linear map $f\colon C_\bullet\to B_{\bullet+r}$ of degree $r$, there is a unique way to extend $f$ to an algebra homomorphism $\SSym\llbracket f\rrbracket$ and to a derivation $\text{Der}(f)$ of degree $r$.

The pairing between $C_\bullet$ and $C^\bullet$ can be promoted to a cap product between $D_\bullet$ and $D^\bullet$. 
\begin{definition}
\source{virasoro-constraints-moduli-sheaves-vertex-algebras}
\notready
\label{Def: pairingofV}
\source{virasoro-constraints-moduli-sheaves-vertex-algebras}
Fix $C_\bullet$ and $C^\bullet$ dual vector spaces and let $\langle-,-\rangle\colon C^\bullet\times C_\bullet\to \BC$ be the pairing. Let $D_\bullet$ and $D^\bullet$ be as in Definition \ref{def: supercommutativefree}.
We define a cap product
$$
\cap\colon C^\bullet\times D_\bullet\longrightarrow D_\bullet
$$
by letting $\nu\cap(-)$ for $\nu\in C^\bullet$ act as a superderivation of degree $-\deg(\nu)$ on $D_\bullet$  restricting to $\langle \nu,-\rangle: C_\bullet\to\BC$. The cap product extends uniquely to
$$
\cap\colon D^\bullet\times D_\bullet\longrightarrow D_\bullet\,,
$$
by requiring that $(\mu \nu) \cap u = \mu\cap(\nu\cap u)$. Notice that this makes $D_{\bullet}$ into a left $D^\bullet$-module.

Starting from a map $f: C_1^\bullet\to C_2^\bullet$, there is a dual map $f^{\dagger}:C_{2,\bullet}\to C_{1,\bullet}$. Constructing algebra homomorphisms and derivations commutes with taking duals: $$\SSym\llbracket f^\dagger\rrbracket = \SSym\llbracket f\rrbracket^ \dagger\,,\qquad \text{Der}(f^\dagger) = \text{Der}(f)^\dagger\,.$$

By composing with the projection $D_\bullet\to \BC$, we recover a non-degenerate pairings $$\langle -,-\rangle: D^\bullet\times D_\bullet\to \BC\,.$$
\end{definition}

An explicit description of the cap product can be obtained after fixing a basis $B\subset C_\bullet$ and observing that 
$$
\nu \cap (-) = \sum_{v\in B}\langle \nu,v\rangle\, \frac{\partial}{\partial v}\,
$$
where $\frac{\partial}{\partial v}$ is the superderivation of degree $-\deg(v)$ acting on the elements of the basis $B$ by $\frac{\partial}{\partial v}(v)=1$ and $\frac{\partial}{\partial v}(w)=0$ for $w\in B\setminus \{v\}$.

\subsection{Descendent algebra}
\label{sec: descendentalgebra}
\source{virasoro-constraints-moduli-sheaves-vertex-algebras}
Let $X$ be a smooth projective variety over $\BC$.

\begin{definition}
\source{virasoro-constraints-moduli-sheaves-vertex-algebras}
\notready\label{def: descendentalgebra}
\source{virasoro-constraints-moduli-sheaves-vertex-algebras}
Let $\text{CH}^X$ denote the infinite dimensional vector space over $\BC$  generated by symbols called \textit{holomorphic descendents} of the form
\[\chh_i(\gamma)\quad\textup{ for }\quad i\geq 0,\, \gamma\in H^\bullet(X)\]
subject to the linearity relations
\[\chh_i(\lambda_1\gamma_1+\lambda_2\gamma_2)=\lambda_1\chh_i(\gamma_1)+\lambda_2\chh_i(\gamma_2)\]
for $\lambda_1, \lambda_2\in \BC$. We define the \textit{cohomological} $\BZ$-grading on $\text{CH}^X$ by 
\begin{equation}
\label{eq: degchigamma}
\source{virasoro-constraints-moduli-sheaves-vertex-algebras}
\deg\chh_i(\gamma) = 2i-p+q \quad \text{ for }\quad \gamma\in H^{p,q}(X)\,.
\end{equation}
Finally, we let $\BD^X$ be the $\BZ$-graded \textit{algebra of holomorphic descendents}
\[\BD^X = \SSym\llbracket\text{CH}^X\rrbracket\,,\]
which is the completion of the supercommutative algebra generated by $\ch_i^\HH(\gamma)$. We will write $\chh_\bullet(\gamma)$ for the element
\[\chh_\bullet(\gamma)=\sum_{i\geq 0}\chh_i(\gamma)\in \BD^X\,.\]
\end{definition}

\begin{remark}
\source{virasoro-constraints-moduli-sheaves-vertex-algebras}
\notready\label{rem: comparingnotation}
\source{virasoro-constraints-moduli-sheaves-vertex-algebras}
This algebra of descendents is very similar to the one introduced in \cite{moop} with two small differences. Firstly, we now take a completion with respect to degree; this makes little difference in practice, but it is important in the comparison between $\BD^X$ and $H^\bullet(\CM_X)$ (cf. Lemma \ref{Lem: isomorphism}) since we follow the standard convention that the total cohomology is a product of the groups in each degree. The second difference is in the grading; in our notation, given $\gamma\in H^{p,q}(X)$, the symbol $\chh_i(\gamma)$ should be thought of as having Hodge degree $(i, i-p+q)$ (recall this from Definition \ref{def: geometricrealization}) so that the geometric realization is degree preserving. The superscript $H$ stands for holomorphic part of the Hodge degree and is used to indicate this degree convention. The original convention appearing in loc. cit. defined the descendents $\ch^{\text{old}}_i(\gamma)$ as \[\ch_i^{\textup{old}}(\gamma)=\chh_{i+p-\dim(X)}(\gamma)\]
for $\gamma\in H^{p,q}(X)$. The reason for introducing holomorphic descendents is to give a natural looking expression for the operator $\mathsf{T}_k$ in Section \ref{sec: Virarosorops}.
\end{remark}

It will sometimes be useful to also consider the shift
\[\ch_{i}(\gamma)\coloneqq \chh_{i+\lfloor\frac{p-q}{2}\rfloor}(\gamma)=\ch^{\textup{old}}_{i-\lceil\frac{\deg \gamma}{2}\rceil+\dim X}(\gamma)\]
so that $\deg\ch_i(\gamma)=2i-|\gamma|$; we recall that $|\gamma|\in \{0,1\}$ is the parity of $\gamma$ as in Section \ref{sec: supercommutative}. 
\begin{definition}
\source{virasoro-constraints-moduli-sheaves-vertex-algebras}
\notready\label{def: topdescendentalgebra}
\source{virasoro-constraints-moduli-sheaves-vertex-algebras}
Let $\alpha\in K_\sst^0(X)$ be a topological type. We define $\tCH^X_\alpha$ to be the graded vector space generated by symbols \[\ch_i(\gamma)\quad\textup{ for }\quad i\in \BZ_{>0},\,\gamma\in H^\bullet(X).\]
We let $\BD^X_\alpha$ be $\SSym\llbracket \tCH^X_\alpha\rrbracket$. The algebra $\BD^X_\alpha$ comes equipped with an algebra homomorphism
$p_\alpha\colon \BD^X\to \BD^X_\alpha$
sending
\[\ch_i^\HH(\gamma)\mapsto\begin{cases}
\ch_{i-\lfloor\frac{p-q}{2}\rfloor}(\gamma) &\textup{ if } \,\deg\ch_i^\HH(\gamma)>0\\
\int_X \gamma\cdot \ch(\alpha) &\textup{ if } \,\deg\ch_i^\HH(\gamma)=0\\
0&\textup{ otherwise.}
\end{cases}\]
\end{definition}

Note that abstractly the algebras $\BD^X_\alpha$ are independent of $\alpha$, but the morphisms $p_\alpha$ depend on $\alpha$ by their behavior on the descendents of degree 0. Let $M_\alpha$ be a moduli space parametrizing sheaves of topological type $\alpha$ with a universal sheaf $\BG$. Then the geometric realization factors through $p_\alpha$:
\begin{equation}
\label{Eq: descentofrealization}
\source{virasoro-constraints-moduli-sheaves-vertex-algebras}
\begin{tikzcd}
\BD^X\arrow[r,"p_{\alpha}"]\arrow[d,swap, "\xi_{\BG}"]& \BD^{X}_\alpha \arrow[dl,"\xi_{\BG}"]\\
H^\ast(M_\alpha)&
\end{tikzcd}
\end{equation} 
where we still denote by $\xi_\BG$ the factoring map, which is defined as
\begin{equation}\label{eq: geometricrealizationalpha}
\source{virasoro-constraints-moduli-sheaves-vertex-algebras}
\xi_{\BG}(\ch_i(\gamma))=p_\ast\left(\ch_{i-\lceil\frac{\deg \gamma}{2}\rceil+\dim X}(\BG)q^\ast \gamma\right)\in H^{2i-|\gamma|}(M_\alpha)\,.
\end{equation}
The map factors since 
$\xi_{\BG}(\chh_i(\gamma))=\int_X \gamma\cdot \ch(\alpha)\in H^0(M_\alpha)$ when $\deg(\chh_i(\gamma))=0$ and the degree of $\xi_{\BG}\big(\chh_i(\gamma)\big)$ is identical to \eqref{eq: degchigamma}. Given $D, D'\in \BD^X$ we say that ``$D=D'$ in $\BD^X_\alpha$'' if $p_\alpha(D)=p_\alpha(D')$.

\begin{example}
\source{virasoro-constraints-moduli-sheaves-vertex-algebras}
\notready
\label{Ex: Tvir}
\source{virasoro-constraints-moduli-sheaves-vertex-algebras}
   One may lift the Chern classes of the virtual tangent bundle of $M_\alpha$ to $\BD^X$ or $\BD^X_{\alpha}$. As a topological K-theory class, the virtual tangent bundle is defined as 
$$
T^\vir M=-Rp_\ast\RHom(\BG, \BG) + \CO_M\,.
$$
 Using $\sum_{t}\gamma_t^L\otimes \gamma_t^R$ to denote the Künneth decomposition of $\Delta_{\ast} \td(X)$, where $\Delta\colon X\to X\times X$ is the diagonal, and applying Grothendieck-Riemann-Roch, one computes that
\begin{align*}
\ch(TM^\vir)=-\xi_\BG\left(\sum_{i,j\geq 0}\sum_t(-1)^{i-p_{t}^L+\dim(X)}\chh_i(\gamma_t^L)\chh_j(\gamma_t^R)\right)+1\,,
\end{align*}
where $\gamma_t^L\in H^{p_t^L, q_t^L}(X)$. The reason for the existence of this lift will become apparent from Lemma \ref{Lem: isomorphism}. The similarity with Virasoro constraints below is a general phenomenon which can be used to guess their correct formulation.
\end{example}



\subsection{Virasoro operators}

\label{sec: Virarosorops}
\source{virasoro-constraints-moduli-sheaves-vertex-algebras}

In this section we define the Virasoro operators 
\[\bL_k\colon \BD^X\to \BD^X\quad\textnormal{for}\quad k\geq -1,\]which produce the Virasoro constraints. These operators have two terms, a derivation term $\bR_k$ and a linear term $\bT_k$ which is quadratic in $\chh_i$. The full Virasoro operators are $\bL_k=\bR_k+\bT_k$, where
\begin{enumerate}
\item 
$\bR_k \colon \BD^X\to \BD^X$ is an even (of degree $2k$) derivation extended from $\bR_k\colon \tCH^X\to\tCH^X$, where it is defined by
\[\bR_k \chh_i(\gamma)\coloneqq\left(\prod_{j=0}^k (i+j)\right)\chh_{i+k}(\gamma)\,.\]
We take the following conventions: the above product is $1$ if $k=-1$ and $\chh_{i+k}(\gamma)=0$ if $i+k<0$. 
\item $\bT_k\colon \BD^X\to \BD^X$ is the operator of multiplication by the element of $\BD^X$ given by
\[\bT_k\coloneqq\sum_{i+j=k}(-1)^{\dim X-p^L}i!j!\chh_i\chh_j(\td(X)).\]
In the formula above, $(-1)^{\dim X-p^L}\chh_i\chh_j(\td(X))$ is defined as follows: let $\Delta\colon X\to X\times X$ be the diagonal map and let 
\[\sum_{t}\gamma_t^L\otimes \gamma_t^R=\Delta_\ast \td(X)\] 
be a Kunneth decomposition of  $\Delta_\ast\td(X)$ such that $\gamma_t^L\in H^{p^L_t, q^L_t}(X)$ for some $p_t^L, q_t^L$. Then 
\[(-1)^{\dim X-p^L}\chh_i\chh_j(\td(X))\coloneqq \sum_{t}(-1)^{\dim X-p_t^L}\chh_i(\gamma_t^L)\chh_j(\gamma_t^R).\]
\end{enumerate}
\begin{remark}
\source{virasoro-constraints-moduli-sheaves-vertex-algebras}
\notready
The operator $\bL_{-1}=\bR_{-1}$ plays a special role and has a particularly nice geometric interpretation in terms of $\BG_m$--gerbes over $M_{\alpha}$ which we describe in Lemma \ref{lem: translationdualr-1}.
\end{remark}

The operators $\{\bL_k\}_{k\geq -1}$ satisfy the Virasoro bracket relations
\[[\bL_k, \bL_\ell]=(\ell-k)\bL_{k+\ell}\in \textup{End}(\BD^X)\]
where the bracket denotes the usual commutator of operators $[X,Y]=X\circ Y-Y\circ X$.
This was noted in \cite{moop}; for a detailed proof see \cite[Proposition 2.10]{bree}. The unusual constant factor $(\ell-k)$, instead of $(k-\ell)$, suggests that there might be another set of more natural Virasoro operators to which $\{\bL_k\}_{k\geq -1}$ are dual. This observation is made into a precise statement in Theorem \ref{thm: duality}. 

\subsection{Weight zero descendents}
\label{sec: invariantdescendents}
\source{virasoro-constraints-moduli-sheaves-vertex-algebras}



One of the issues that arise when dealing with descendent invariants is that a priori they depend on a choice of universal sheaf $\BG$ on $M\times X$. Given a universal sheaf $\BG$, the possible universal sheaves are of the form
$\BG'=\BG\otimes p^\ast L$ where $L$ is a line bundle on $M$.
\begin{lemma}
\source{virasoro-constraints-moduli-sheaves-vertex-algebras}
\notready
\label{lem: changeuniversalsheaf}
\source{virasoro-constraints-moduli-sheaves-vertex-algebras}
Let $\bE\colon \BD^X\to \BD^X\llbracket\zeta\rrbracket$ be the algebra homomorphism defined by
$$
\bE = e^{\zeta \bR_{-1}}\,,
$$
then given two universal sheaves $\BG$ and $\BG'=\BG\otimes p^\ast L$, the geometric realizations with respect to the two are related by
\[\xi_{\BG'}=\xi_{\BG}\circ \bE|_{\zeta = c_1(L)} \,.\]
\end{lemma}
\begin{proof}
We may write 
\[\bE\big(\chh_i(\gamma)\big) =\sum_{n=0}^i \frac{\zeta^n}{n!}\chh_{i-n}(\gamma)= [e^\zeta \chh_\bullet(\gamma)]_i\] 
which after comparing to
\[\xi_{\BG'}(\chh_\bullet(\gamma))=p_\ast\left(\ch(\BG \otimes p^\ast L)q^\ast \gamma\right)=e^{c_1(L)}\xi_{\BG}(\chh_\bullet(\gamma)).\qedhere\]
yields the result.
\end{proof}
In particular, it follows that if $D$ is such that $\bR_{-1}(D)=0$ then $\xi_\BG(D)=\xi_{\BG'}(D)$. This leads to the definition of the algebra of weight 0 descendents. Its geometric interpretation is summarized in Lemma \ref{Lem: DXinv} and is related to taking rigidification of moduli stacks. Roughly speaking, these classify families of sheaves on $S\times X$ 
up to twisting by line bundles on $L\to S$ and the above discussion formulates the precise interaction between the twisting and the descendents. 



\begin{definition}
\source{virasoro-constraints-moduli-sheaves-vertex-algebras}
\notready\label{def: invariantdescendents}
\source{virasoro-constraints-moduli-sheaves-vertex-algebras}
For a topological type $\alpha\in K_\sst^0(X)$, we will also denote by 
\[\bR_{-1}\colon \BD^X_{ \alpha}\to \BD^X_{ \alpha}\] the derivation defined on generators by
$\bR_{-1}\ch_i(\gamma)=\ch_{i-1}(\gamma)$,
where $\ch_0(\gamma)$ is interpreted as $\int_X \gamma\cdot \ch(\alpha)$.
We then define
\begin{align*}\BD^X_{\inv}&=\{D\in \BD^X\colon \bR_{-1}(D)=0\}\,,\\
\BD^X_{\inv, \alpha}&=\{D\in \BD^X_{\alpha}\colon \bR_{-1}(D)=0\}\,.
\end{align*} 
\end{definition}

For any weight 0 descendent $D\in \BD^X_{\inv}$, by Lemma \ref{lem: changeuniversalsheaf}, the geometric realization $\xi_{\BG}(D)$ does not depend on the choice of $\BG$. Thus, when $D\in \BD^X_{\inv}$ we will often omit specifying the realization map and write
\[\int_{[M]^\vir} D=\int_{[M]^\vir}\xi_{\BG}(D)\]
for \emph{any} choice of universal sheaf $\BG$. 
The morphism $\BD^X_{\inv}\to H^\bullet(M)$ is defined even without assuming the existence of any universal sheaf $\BG$. This fact can be proven using Lemma \ref{Lem: DXinv}; if $M=M_\alpha$ is a moduli space of topological type $\alpha$, then it admits an open embedding $\iota\colon M\hookrightarrow \CM_\alpha^\rig$ and thus we get a map
\[\begin{tikzcd}
\BD^X_{\inv}\arrow[r, "p_\alpha"]&
\BD^X_{\inv, \alpha}\arrow[r, "\sim"]&
H^\bullet(\CM_\alpha^\rig)\arrow[r, "\iota^\ast"]&
H^\bullet(M).
\end{tikzcd}\]

\begin{example}
\source{virasoro-constraints-moduli-sheaves-vertex-algebras}
\notready\label{Ex: DXinvdescendents}
\source{virasoro-constraints-moduli-sheaves-vertex-algebras}
Given $\gamma_1,\gamma_2\in H^\bullet(X)$ we have
\[\chh_1(\gamma_1)\chh_0(\gamma_2)-\chh_0(\gamma_1)\chh_1(\gamma_2)\in \BD^X_{\inv}.\]
One can also check that the lift of $\ch(T^\vir M)$ to $\BD^X$ is in $\BD^X_{\inv}$ using the expression in Example \ref{Ex: Tvir}. A geometric reason for this is going to be given in Example \ref{ex: invdescendents}.
\end{example}



\subsection{Virasoro constraints for sheaves}
\label{sec: virasorosheaves}
\source{virasoro-constraints-moduli-sheaves-vertex-algebras}

To formulate the Virasoro constraints for moduli of sheaves, without a canonical choice of universal sheaf, we must produce relations among weight 0 descendents. This is achieved by combining the Virasoro operators previously introduced in the way which we now describe:
\begin{definition}
\source{virasoro-constraints-moduli-sheaves-vertex-algebras}
\notready\label{def: Linv}
\source{virasoro-constraints-moduli-sheaves-vertex-algebras}
The weight 0 Virasoro operator $\Linv \colon \BD^X\to \BD^X$ is defined by
\[\Linv=\sum_{j\geq -1}\frac{(-1)^j}{(j+1)!}\bL_j \bR_{-1}^{j+1}.\]
\end{definition}

The operator $\Linv$ maps $\BD^X$ to $\BD^X_{\inv}$. Indeed, using that $[\bR_{-1}, \bL_j]=(j+1)\bL_{j-1}$ we find
\begin{align*}
    \bR_{-1}\circ \Linv &=\sum_{j\geq -1}\frac{(-1)^j}{(j+1)!}\bL_j \bR_{-1}^{j+2}+\sum_{j\geq -1}\frac{(-1)^j}{(j+1)!}(j+1)\bL_{j-1} \bR_{-1}^{j+1}=0\,.
\end{align*}

In particular, for any $D\in \BD^X$ the integral $\int_M \xi_{\BG}(\Linv(D))$ does not depend on the universal sheaf $\BG$ so we omit the realization homomorphism $\xi_{\BG}$ from the notation. 

\begin{conjecture}
\source{virasoro-constraints-moduli-sheaves-vertex-algebras}
\notready\label{conj: invVirasoro}
\source{virasoro-constraints-moduli-sheaves-vertex-algebras}
Let $M$ be a moduli of sheaves as in Section \ref{sec: modulisheavespairs}. Then
\[\int_{[M]^\vir} \Linv(D)=0\quad \textup{ for any }D\in \BD^X.\]
\end{conjecture}

This formulation is different from the ones which appear in previous works, namely \cite{moop, moreira, bree}. There, Virasoro constraints are formulated using a choice of universal sheaf that is natural in each of the moduli spaces considered. 

\begin{definition}
\source{virasoro-constraints-moduli-sheaves-vertex-algebras}
\notready
Let $M=M_\alpha$ be a moduli space of sheaves of topological type $\alpha$ and let $\delta\in H^\bullet(X, \BZ)$ be an algebraic class such that $\int_X  \delta\cdot \ch(\alpha) \neq 0$. We say that a universal sheaf $\BG$ is $\delta$-normalized if
\begin{equation}\label{Eq:deltanorm}\xi_\BG(\chh_1(\delta))=0.\end{equation}
\source{virasoro-constraints-moduli-sheaves-vertex-algebras}
\end{definition}

Given $\delta$ as in the definition, we define the operators
\[\bL^\delta_k=\bR_k+\bT_k+\bS_k^\delta,\quad k\geq -1\,,\]
where
\[\bS_k^\delta=-\frac{(k+1)!}{\int_X \delta\cdot \ch(\alpha)}\bR_{-1}\circ \chh_{k+1}(\delta).\]

\begin{remark}
\source{virasoro-constraints-moduli-sheaves-vertex-algebras}
\notready\label{rem: rationaldeltanormalized}
\source{virasoro-constraints-moduli-sheaves-vertex-algebras}
If $\delta$ is such that $\int_X  \delta\cdot \ch(\alpha)\neq 0$ then a $\delta$-normalized universal sheaf always exists (and is unique) as an element of the rational $K$-theory of $M\times X$. Precisely, for any universal sheaf $\BG$ we have
\[\BG_\delta=\BG\otimes e^{-\xi_\BG(\chh_1(\delta))/\int_X \delta\cdot \ch(\alpha)}\]
in the rational $K$-theory of $M\times X$ that can be thought of as the unique $\delta$-normalized universal sheaf; here we use $e^{c}$ to denote a rational line bundle with first Chern class equal to the algebraic class $c\in H^2(M, \BQ)$. 

The geometric realization with respect to $\BG_\delta$ is given by $\xi_{\BG_\delta}=\xi_{\BG}\circ \eta$ where $\eta\colon \BD^X\to \BD^X_{\inv,\alpha}$ is 
\[\eta=\sum_{j\geq 0}\frac{1}{j!}\left(-\frac{\chh_1(\delta)}{\int_X \delta\cdot \ch(\alpha)}\right)^j\bR_{-1}^j\,.\]
Thus, we can talk about the geometric realization with respect to a $\delta$-normalized sheaf even if such a sheaf does not exist in the usual sense. Conjecture \ref{conj: normalizedVirasoro} still makes sense in this setting and the proof of Proposition \ref{prop: normalizedvsinv} goes through as well.
\end{remark}

\begin{conjecture}
\source{virasoro-constraints-moduli-sheaves-vertex-algebras}
\notready\label{conj: normalizedVirasoro}
\source{virasoro-constraints-moduli-sheaves-vertex-algebras}
Let $M=M_\alpha$ be a moduli of sheaves as in Section \ref{sec: modulisheavespairs} and let $\BG$ be a $\delta$-normalized universal sheaf. Then
\[\int_{[M]^\vir} \xi_\BG(\bL_k^\delta(D))=0  \quad \textup{for any }k\geq -1, \, D\in \BD^X.\]
\end{conjecture}

\begin{proposition}
\source{virasoro-constraints-moduli-sheaves-vertex-algebras}
\notready\label{prop: normalizedvsinv}
\source{virasoro-constraints-moduli-sheaves-vertex-algebras}
Conjectures \ref{conj: invVirasoro} and \ref{conj: normalizedVirasoro} are equivalent.
\end{proposition}
\begin{proof}
We begin by observing that we have the identity
\begin{equation}\label{eq: invLnormalized}
\source{virasoro-constraints-moduli-sheaves-vertex-algebras}
\Linv=\sum_{j\geq -1}\frac{(-1)^j}{(j+1)!}\bL_j^\delta \bR_{-1}^{j+1}\,,
\end{equation}
that follows from
\begin{align*}
\sum_{j\geq -1}&\frac{(-1)^j}{(j+1)!}\bS_j^\delta \bR_{-1}^{j+1}\\
&=-\frac{1}{\int_X \delta\cdot \ch(\alpha)}\sum_{j\geq -1}(-1)^j\left( \chh_j(\delta) \bR_{-1}^{j+1}+\chh_{j+1}(\delta)\bR_{-1}^{j+2}\right)=0.
\end{align*}
By \eqref{eq: invLnormalized}, Conjecture \ref{conj: normalizedVirasoro} clearly implies \ref{conj: invVirasoro}.

For the reverse implication we use (backward) induction on $k$. For every $k>\mathrm{virdim}(M)$ the statement of Conjecture $\ref{conj: normalizedVirasoro}$ is clear by degree reasons. Assume now that the result holds for every $k'>k$, and let $F=\chh_1(\delta)\in \BD^X$ satisfying $\xi_\BG(F)~=~0$ by \eqref{Eq:deltanorm}. The weight 0 Virasoro operator applied to $F^{k+1}D$ gives
\begin{equation}\label{eq: invariantimpliesnormalized}
\source{virasoro-constraints-moduli-sheaves-vertex-algebras}
0=\int_{[M]^\vir} \xi_{\BG}\left(\Linv(F^{k+1}D)\right)=\sum_{j\geq -1}\frac{(-1)^j}{(j+1)!}\int_{[M]^\vir}\xi_{\BG}\left(\bL_j^\delta \bR_{-1}^{j+1}(F^{k+1}D)\right).
\end{equation}
Since $\bR_{-1}$ is a derivation and $\bR_{-1}(F)=\chh_0(\delta)=\int_X \delta\cdot \ch(\alpha)$ in $\BD^{X}_\alpha$, we have
\begin{align}\label{eq: invariantimpliesnormalized2}
\source{virasoro-constraints-moduli-sheaves-vertex-algebras}
\nonumber \bR_{-1}^{j+1}(F^{k+1}D)&=\sum_{s=0}^{j+1} \binom{j+1}{s}\bR_{-1}^s(F^{k+1})\bR_{-1}^{j+1-s}(D)\\
&=\sum_{s=0}^{\min(k+1, j+1)}\binom{j+1}{s}\frac{(k+1)!}{(k+1-s)!}r^s F^{k+1-s}R_{-1}^{j+1-s}(D)
\end{align}
in $\BD^{X}_{\alpha}$, where we denote $r=\int_X \delta\cdot \ch(\alpha)$. Note also that the operators $\bL_j^\delta$ satisfy the property that $\bL_{j}^\delta(F D)=F \,\bL_{j}^\delta(D)$ in $\BD^{X}_\alpha$ for every $D$. Since by \eqref{Eq:deltanorm} the geometric realization of $F$ with respect to $\BG$ is zero, the only terms of \eqref{eq: invariantimpliesnormalized2} contributing to \eqref{eq: invariantimpliesnormalized} are the ones with $s=k+1$, thus  \eqref{eq: invariantimpliesnormalized} becomes
\[0=\sum_{j\geq k}\frac{(-1)^j}{(j-k)!} r^{k+1}\int_{[M]^\vir}\xi_{\BG}\left(\bL_j^\delta \bR_{-1}^{j-k}(D)\right)\]
By the induction hypothesis, all the terms with $j>k$ vanish and thus the term with $j=k$ also vanishes, showing that 
\[0=\int_{[M]^\vir} \xi_{\BG}(\bL_k^\delta(D)),\]
which concludes the induction step.\qedhere
\end{proof}

\begin{remark}
\source{virasoro-constraints-moduli-sheaves-vertex-algebras}
\notready\label{rem: normalizedpt}
\source{virasoro-constraints-moduli-sheaves-vertex-algebras}
For surfaces or 3-folds $X$ with $H^i(\CO_X)=0$ for $i>0$, the Hilbert scheme of points and the moduli of stable pairs, respectively, are moduli of sheaves in the sense of Section \ref{sec: modulisheavespairs}. The canonical universal sheaves in each of these cases are precisely the $\pt$-normalized universal sheaves and the formulations in \cite{moop, moreira} coincide with Conjecture \ref{conj: normalizedVirasoro}. In \cite{bree}, the author considers moduli spaces of torsion free stable sheaves on surfaces with $h^{1,0}=h^{2,0}=0$; the sheaves in such moduli spaces have fixed determinant $\Delta$. The author uses a geometric realization with respect to the sheaf
$\BG\otimes \det(\BG)^{-1/r}\,.$
This is not a universal sheaf for $M$ in the sense we use in this paper since its restriction to $\{G\}\times X$ is $G\otimes \det(G)^{-1/r}=G\otimes \Delta^{-1/r}$ rather than $G$; however,
\[\BG\otimes \det(\BG)^{-1/r}\otimes q^\ast\Delta^{1/r}\]
recovers the $\pt$-normalized universal sheaf of Remark \ref{rem: rationaldeltanormalized}. The equivalence between van Bree's formulation of Virasoro for the moduli of stable sheaves of positive rank and Conjecture \ref{conj: normalizedVirasoro} is explained by Lemma \ref{lem: twist}.
\end{remark}

\subsection{Virasoro constraints for pairs}
\label{sec: virasoropairs}
\source{virasoro-constraints-moduli-sheaves-vertex-algebras}

Since for moduli spaces of pairs we have a uniquely defined universal object $q^\ast V\to \BF$, there is no need to use the weight 0 Virasoro operator for pairs. We need, however, to slightly modify the operators to account for the different obstruction theory.

Let $V$ be a fixed sheaf on $X$ and let $P$ as in  Section \ref{sec: modulisheavespairs} be a moduli space of pairs parametrizing sheaves $F$ together with a map $V\to F$. The moduli $P$ comes equipped with a (unique) universal sheaf $\BF$ on $P\times X$ and a universal map $q^\ast V\to \BF$. We conjecture that descendent invariants obtained by integration on moduli of pairs are constrained by Virasoro operators which are similar to the ones introduced in Section \ref{sec: Virarosorops}. Define the operators 
\[\bL_k^V\colon \BD^X\to \BD^X\]
by $\bL_k^V=\bR_k+\bT_k^V$ where
\[\bT_k^V=\bT_k-k!\chh_k\left(\ch(V)^\vee\td(X)\right).\]
The operator $\bL_k^V$ can be described in an alternative way that should make its definition more natural and that will be closer to the vertex algebra language that we will introduce later. Let 
\[\Dpa=\BD^X\otimes\BD^X\]
be the algebra of pair descendents. We denote the generators of the first copy of $\BD^X$ by  $\ch_i^{\HH, \CV}(\gamma)$ and the generators of the second copy by  $\ch_i^{\HH, \CF}(\gamma)$. Given the universal pair $q^\ast V\to \BF$ on $P\times X$, we have a geometric realization map
\[\xi_{(q^\ast V,\BF)}\colon \Dpa\to H^\bullet(P)\]
that is defined by
\begin{align*}\xi_{(q^\ast V, \BF)}\left(\ch_i^{\HH, \CV}(\gamma)\right)&=\xi_{q^\ast V}\left(\chh_i(\gamma)\right)=\begin{cases}\int_X \gamma\cdot \ch(V)&\textup{ if }i=0\\
0 & \textup{ otherwise}\end{cases}\\
\xi_{( q^\ast V, \BF)}\left(\ch_i^{\HH, \CF}(\gamma)\right)&=\xi_{\BF}\left(\chh_i(\gamma)\right).
\end{align*}
This geometric realization map factors through
\begin{center}
\begin{tikzcd}
\Dpa\arrow[d, swap, "\xi_V"]\arrow[rd, "\xi_{(q^\ast V,\BF)}"]& \\
\BD^X\arrow[r, "\xi_{\BF}"]&
H^\bullet(P)
\end{tikzcd}
\end{center}
where $\xi_{V}$ is defined to send 
\[\ch_i^{\HH, \CV}(\gamma)\mapsto \delta_i \int_X \gamma\cdot \ch(V)\textup{ and }\ch_i^{\HH, \CF}(\gamma)\mapsto \ch_i^{\HH}(\gamma)\,;\]
here $\delta_i=\delta_{i,0}$ is equal to $1$ if $i=0$ and $0$ otherwise.
We define the pair Virasoro operators $\bL_k^{\pa}\colon \Dpa\to \Dpa$, for $k\geq -1$, as the sum $\bR_k^{\pa}+\bT_k^{\pa}$ where
\begin{enumerate}
\item $\bR_k^{\pa}$ is a derivation defined on generators in the same way as $\bR_k$; in other words,
\[\bR_k^{\pa}=\bR_k\otimes \id+\id\otimes \bR_k.\]
\item $\bT_k^{\pa}$ is the operator of multiplication by the element
\[\bT_k^{\pa}=\sum_{i+j=k} (-1)^{\dim X- p^L}i!j!\ch_i^{\HH, \CF-\CV}\ch_j^{\HH, \CF}(\td(X))\in \Dpa\]
where
\[\ch_i^{\HH, \CF-\CV}\coloneqq \ch_i^{\HH, \CF}-\ch_i^{\HH, \CV}.\]
\end{enumerate}


The definition of $\bT_k^{\pa}$ suggests an intimate relation between the linear part of the Virasoro operators and the obstruction theory of the moduli spaces we consider; recall that the deformation-obstruction theory of the pair moduli space $P$ at $V\to F$ is governed by
\[\RHom\left([V\to F], F\right)\]
and use Example \ref{Ex: Tvir}.

The operators $\bL_k^V$ are obtained from $\bL_k^{\pa}$ by the following commutative diagram:
\begin{center}
\begin{tikzcd}
\Dpa\arrow[r, "\bL_k^{\pa}"]\arrow[d, "\xi_V"]&
\Dpa\arrow[d, "\xi_V"]\\
\BD^X\arrow[r, "\bL_k^V"]&
\BD^X.
\end{tikzcd}
\end{center}
This holds due to the identity
\begin{align*}
\sum_{t}(-1)^{\dim X-p_t^L}\left(\int_X \gamma_t^L \cdot \ch(V) \right)\chh_k(\gamma_t^R)&=\sum_t \left(\int_X \gamma_t^L \cdot \ch(V)^\vee \right)\chh_k(\gamma_t^R)\\
&=\chh_k(\ch(V)^\vee \td(X)).
\end{align*}

% define $\ch_k^{H, \BF-V}(\gamma)$ by \wcomment{I think it would be better to use $\ch_k^{H, \BF-V}$ exclusively than $\ch_k^{H,V-\BF}$. This for instance remove $-$ sign in $T_k^V$ that makes it closer to previous $T_k$.}
%\[\ch_k^{H, V-\BF}(\gamma)=\delta_k \int_X \gamma\cdot \ch(V)-\chh_k(\gamma).\]
%The notation is explained by the fact that
%\[\xi_{\BF}\left(\ch_\bullet^{H, V-\BF}(\gamma)%\right)=p_\ast\left(\ch_\bullet(q^\ast V-\BF)q^\ast% \gamma\right).\]
%Then $T_k^V$ can be expressed as
%\begin{equation}\label{eq: TkV}
%T_k^V=-\sum_{a+b=k}\sum_i (-1)^{\dim X- p_i^L}a!b!\ch_a^{H, V-\BF}\chh_b(\td(X))
%\end{equation}

%The expression \eqref{eq: TkV} suggests an intimate relation between the linear part of the Virasoro operators and the obstruction theory of the moduli spaces we consider; recall that the deformation-obstruction theory of the pair moduli space $P$ at $V\to F$ is governed by
%\[\RHom\left([V\to F], F\right).\]

It can be shown that the operators $\{\bL_k^V\}_{k\geq -1}$ and $\{\bL_k^{\pa}\}_{k\geq -1}$ satisfy the Virasoro bracket relations.

\begin{conjecture}
\source{virasoro-constraints-moduli-sheaves-vertex-algebras}
\notready\label{conj: pairvirasoro}
\source{virasoro-constraints-moduli-sheaves-vertex-algebras}
Let $P$ be a moduli of pairs as in Section \ref{sec: modulisheavespairs} with universal pair $q^\ast V\to \BF$ on $P\times X$. Then
\[\int_{[P]^\vir}\xi_\BF\big(\bL_k^V(D)\big)=0\quad \textup{ for any }k\geq 0, \, D\in \BD^X.\]
\end{conjecture}
Equivalently, the pair Virasoro constraints can be formulated as
\[\int_{[P]^\vir} \xi_{(q^\ast V,\BF)}\big(\bL_k^{\pa}(D)\big)=0\quad \textup{ for any }k\geq 0, \, D\in \Dpa.\]

\subsection{Invariance under twist}
\label{sec: invariancetwist}
\source{virasoro-constraints-moduli-sheaves-vertex-algebras}


Suppose that $M=M_\alpha$ is the moduli space of slope semistable sheaves with respect to a polarization $H$ with topological type $\alpha$. Then, the moduli space $M'=M_{\alpha(H)}$ is isomorphic to $M$ by sending a sheaf $[F]\in M$ to $[F']=[F\otimes H]$. As expected, the Virasoro constraints on $M$ and $M'$ are equivalent as we now proceed to verify. 

Suppose that $\BG$ is a universal sheaf on $M\times X$. The universal sheaf on $M'\times X$ is identified with $\BG'=\BG\otimes q^\ast H$ via the isomorphism $M'\times X\cong M\times X$. Define an algebra isomorphism $\mathsf F\colon \BD^X\to \BD^X$ by
\[\bF(\chh_i(\gamma))=\chh_i(e^{c_1(H)} \gamma)\,.\]
Then the following diagram commutes:
\begin{center}
\begin{tikzcd}
\BD^X\arrow[r, "\xi_{\BG'}"] \arrow[d, "\bF"]& H^\ast(M')\arrow[d, "\sim" {rotate=90, below}]\\
\BD^X\arrow[r, "\xi_{\BG}"] & H^\ast(M)
\end{tikzcd}
\end{center}


\begin{lemma}
\source{virasoro-constraints-moduli-sheaves-vertex-algebras}
\notready\label{lem: twist}
\source{virasoro-constraints-moduli-sheaves-vertex-algebras}
The isomorphism $\bF$ commutes with the Virasoro operators, i.e.\[\bL_k\circ \bF=\bF\circ \bL_k\textup{ for $k\geq -1\quad$   and }\quad\Linv\circ \bF=\bF\circ \Linv.\]
Thus, Conjecture \ref{conj: invVirasoro} holds for $M$ if and only if it holds for $M'$.	
\end{lemma}
\begin{proof}
It is enough to show that $\bL_k$ and $\bF$ commute. Commutativity with $\Linv$ follows immediately from its definition; the equivalence of Virasoro constraints follows from the commutativity and the diagram before the Lemma.

The commutativity with the derivation part, i.e., $\bR_k\circ \bF=\bF\circ \bR_k$, is straightforward. To show commutativity with the operator $\bT_k$ we need $\bF(\bT_k)=\bT_k$ (recall that $\bT_k$ denotes both an element in $\BD^X$ and the operator which is multiplication by that element). 
\begin{align*}\bF(\bT_k)&=\sum_{i+j=k}\sum_t(-1)^{\dim X-p^L_t}i!j!\chh_i(e^{c_1(H)}\gamma_t^L)\chh_j(e^{c_1(H)}\gamma_t^R)\\
&=\sum_{i+j=k}\sum_t \sum_{a, b\geq 0}(-1)^{\dim X-p^L_t}i!j!\frac{1}{a!b!}\chh_i(c_1(H)^a\gamma_t^L)\chh_j(c_1(H)^b\gamma_t^R)\\
&=\sum_{i+j=k} \sum_{a, b\geq 0}(-1)^{\dim X-p^L-a}i!j!\frac{1}{a!b!}\chh_i\chh_j(c_1(H)^{a+b}\td(X))\\
&=\sum_{i+j=k} (-1)^{\dim X-p^L}i!j!\chh_i\chh_j(\td(X))=\bT_k.
\end{align*}
The last line uses the fact that for fixed $c>0$ the sum $\sum\limits_{a+b=c}\frac{(-1)^a}{a!b!}$ vanishes.\qedhere
\end{proof}



\subsection{Variants for the fixed determinant theory}
\label{sec: fixeddet}
\source{virasoro-constraints-moduli-sheaves-vertex-algebras}

As we pointed out in Remark \ref{rem: tracelessobs}, this paper is mostly concerned with moduli spaces of sheaves without fixed determinant; in other words, our obstruction theories use full $\textup{Ext}$ groups instead of traceless $\textup{Ext}$ groups. This contrasts with the situation for Pandharipande-Thomas stable pairs or Hilbert schemes of points studied in \cite{moreira}. There, we see a new term appearing in the Virasoro operators which we now recall.

\begin{definition}
\source{virasoro-constraints-moduli-sheaves-vertex-algebras}
\notready
Given a class $\gamma\in H^\bullet(X)$ denote by $\bR_{-1}[\gamma]$ the superderivation (of degree $\deg(\gamma)-2)$ acting on generators by
\[\bR_{-1}[\gamma]\chh_i(\gamma')=\chh_{i-1}(\gamma \gamma').\]
For $k\geq -1$ we define the operator $\bS_k\colon \BD^X\to \BD^X$ by
\[\bS_k=(k+1)!\sum_{p_t^L=0}\bR_{-1}[\gamma_t^L]\circ\chh_{k+1}(\gamma_t^R)\]
where the sum runs over the terms $\gamma_t^L\otimes\gamma_t^R$ in the K\"unneth decomposition of $\Delta_\ast 1$ such that $p_t^L=0$.
\end{definition}
\begin{remark}
\source{virasoro-constraints-moduli-sheaves-vertex-algebras}
\notready
The part of $\bS_k$ corresponding to the K\"unneth component $1\otimes \pt$ is equal to $-r\bS_k^\pt$ where $r=\ch_0(\alpha)=\rk(\alpha)$. When $h^{0, q}(X)=0$ for $q>0$ the appearance of the operator $\bS_k$ is already explained in Remark \ref{rem: normalizedpt} as being related to the $\pt$-normalized universal sheaf. 
\end{remark}

We now proceed to explain the appearance of the operator $\bS_k$ in the more general case in which $h^{0, 1}\neq 0$ but $h^{0, q}=0$ for $q>1$. Let $\alpha$ be such that $r=\rk(\alpha)>0$ and $M=M_\alpha$ be a moduli space parametrizing semistable sheaves with topological type~$\alpha$. Let $\Delta\in \Pic(X)$ be a fixed line bundle  such that $c_1(\Delta)=c_1(\alpha)$. We let $M_\Delta\subseteq M$ be the moduli space of sheaves on $F\in M$ with fixed determiant $\det(F)=\Delta$; i.e., $M_\Delta$ is the pullback
\begin{equation}\label{eq: cartesianpic}
\source{virasoro-constraints-moduli-sheaves-vertex-algebras}
\begin{tikzcd}
M_\Delta\arrow[r, hookrightarrow]\arrow[d]& M\arrow[d, "\det"]\\
\{\Delta\} \arrow[r, hookrightarrow]&
\Pic^{c_1(\alpha)}(X).
\end{tikzcd}
\end{equation}
The moduli space $M_\Delta$ has a 2-term obstruction theory given by 
\begin{align*}
\textup{Tan}&=\Ext^1(G,G)_0\\
\textup{Obs}&=\Ext^2(G,G)_0=\Ext^2(G, G)\\
0&=\Ext^{>2}(G,G).
\end{align*}
Given such a moduli space, there is a unique universal sheaf (possibly rational, in the same sense of Remark \ref{rem: rationaldeltanormalized}) $\BG_\Delta$ on $M_\Delta\times X$ such that $\det(\BG_\Delta)=q^\ast \Delta$. 
We suppose also that the Jacobian $\Pic^0(X)$ acts on $M$ in the natural way; that is,
\[[F]\in M\Rightarrow [F\otimes L]\in M\]
for $L\in \Pic^0(X)$. The two main examples to keep in mind are
\begin{enumerate}
\item $M=M_C(r, d)$ a moduli of stable bundles on a curve $C$ and $\Delta$ a line bundle of degree $d$;
\item $M=M_S^H(r, c_1,c_2)$ a moduli of stable sheaves on a surface $S$ with $p_g(S)=0$ and $\Delta$ a line bundle with $c_1(\Delta)=c_1$. Note that if we take $r=1, c_1=0$ and $\Delta=\CO_S$ we recover the Hilbert scheme of points on $S$
\[M_S^H(1, \CO_S, n)= S^{[n]}\,.\]
\end{enumerate}



\begin{proposition}
\source{virasoro-constraints-moduli-sheaves-vertex-algebras}
\notready\label{prop: fixeddet}
\source{virasoro-constraints-moduli-sheaves-vertex-algebras}
Suppose that $X$ is such that $h^{0, q}=0$ for $q>1$ and $M, M_\Delta$ are as described before. Suppose also that the Virasoro constraints (Conjecture \ref{conj: invVirasoro}) hold for $M$.  Then, we have
\[\int_{[M_\Delta]^\vir}\xi_{\BG_\Delta} \big(\mathcal L_k(D)\big)=0\quad \textup{ for any }k\geq -1, \, D\in \BD^X\]
where 
\[\mathcal L_k=\bL_k-\frac{1}{r}\bS_k\,.\footnote{The $-$ sign in the $\bS_k$ operator does not appear in \cite{moreira} due to the fact that in \cite{moreira} the geometric realization is taken with respect to $-\BG_\Delta$ instead of $\BG_\Delta$.}\]
\end{proposition}
\begin{proof}
We use the tensoring morphism
\[\pi\colon \widetilde M=M_\Delta\times \Pic^0(X)\to M,\quad (F,L)\mapsto F\otimes L.\]
This morphism is a part of the Cartesian diagram 
\begin{center}
\begin{tikzcd}
M_\Delta\times \Pic^0(X) \arrow[r, "\pi"] \arrow[d, "i\times \id"', hook] & M \arrow[d, "{(\id,\det)}", hook] \\
M\times \Pic^0(X) \arrow[r, "f"]                       & M\times \Pic^{c_1(\alpha)}(X)                      
\end{tikzcd}
\end{center}
where $f(G,L)\coloneqq(G\otimes L,\Delta\otimes L^r)$ with $r=\rk(\alpha)$. The morphism $f$ is essentially the multiplication by $r$ map of the abelian variety $\Pic^0(X)$. This implies that $\pi$ is an \'etale morphism of degree $r^{2g}$ where $g=h^{0,1}(X)=\dim(\Pic^0(X))$. 

On the other hand, $f$ respects the obstruction theory because the obstruction theory of $M$ is invariant under tensoring by line bundles. This implies the compatibility between virtual classes
\[\pi^*[M]^\vir=[M_\Delta]^\vir\times [\Pic^0(X)]=:[\widetilde M]^\vir
\]
and $\pi_*[\widetilde M]^\vir=r^{2g}\cdot[M]^\vir$.

We now apply the above discussion to translate the Virasoro constraints of $M$ to that of $M_\Delta$. Let $\CP$ be the Poincaré bundle on $\Pic^0(X)\times X$ (i.e., the $\pt$-normalized universal bundle of $\Pic^0(X)$) and let $\BG$ be the $\pt$-normalized universal sheaf on $M\times X$. Since $\widetilde{\BG}=\BG_\Delta\boxtimes \CP$ is also $\pt$-normalized it follows that $(\pi\times \id)^\ast \BG=\widetilde \BG$ (in rational $K$-theory) and thus $\pi^\ast\circ \xi_{\BG}=\xi_{\widetilde \BG}$. Using the compatibilities between virtual classes and Proposition \ref{prop: normalizedvsinv} we get
\begin{equation}\label{eq: virasoromtilde}
\source{virasoro-constraints-moduli-sheaves-vertex-algebras}
\int_{[\widetilde M]^\vir}\xi_{\widetilde \BG}\big(\big(\bL_k+\bS_k^\pt\big)(D)\big)=0\,
\end{equation}
where $$\bS_k^\pt=-\frac{1}{r}\bR_{-1}\ch_{k+1}(\pt)\,.$$
%Let 
%\[\pi\colon \widetilde M=M_\Delta\times \Pic^0(X)\to M\]
%be the map sending $(F, L)$ to $F\otimes L$. Let $\CP$ be the Poincaré bundle on $\Pic^0(X)\times X$ (i.e., the $\pt$-normalized universal bundle of $\Pic^0(X)$) and let $\BG$ be the $\pt$-normalized universal sheaf on $M\times X$. Since $\widetilde{\BG}=\BG_\Delta\boxtimes \CP$ is also $\pt$-normalized it follows that $(\pi\times \id)^\ast \BG=\widetilde \BG$ (in rational $K$-theory) and thus $\pi^\ast\circ \xi_{\BG}=\xi_{\widetilde \BG}$. Using that  
%\[\pi^*[M]^\vir=[\widetilde M]^\vir\coloneqq [M_\Delta]^\vir\times[\Pic^0(X)]\]
%and Proposition \ref{prop: normalizedvsinv} we get
%\begin{equation}\label{eq: virasoromtilde}
%\int_{[\widetilde M]^\vir}\xi_{\widetilde \BG}\big(\big(\bL_k+\bS_k^\pt\big)(D)\big)=0\,
%\end{equation}
%by the push-pull formula and the fact that $\pi_\ast 1=\deg(\pi)$, where $$\bS_k^\pt=-\frac{1}{r}\bR_{-1}\ch_{k+1}(\pt)\,.$$
Recall that $g=h^{0,1}(X)=h^{1,0}(X)$ and let $\{e_1, \ldots, e_g\}\subseteq H^{0,1}(X)$ and $\{f_1, \ldots, f_g\}\subseteq H^{1,0}(X)$ be basis; let $\{\hat e_1, \ldots, \hat e_g\}\subseteq H^{m, m-1}(X)$, $\{\hat f_1, \ldots, \hat f_g\}\subseteq H^{m-1, m}(X)$ be their dual basis, meaning that 
\[\int_X e_i \hat e_j=\delta_{ij}=\int_X f_i \hat f_j\,.\]
Let 
\[\tau_j=\xi_{\CP}(\chh_1(\hat e_j))\in H^{1,0}(\Pic^0(X))\quad \textup{and}\quad \rho_j=\xi_{\CP}(\chh_0(\hat f_j))\in H^{0,1}(\Pic^0(X))\]
so that
\[c_1(\CP)=\sum_{j=1}^g \tau_j\otimes e_j+\sum_{j=1}^g \rho_j\otimes f_j\in H^2(\Pic^0(X)\times X).\]
The Jacobian is topologically a real torus of dimension $2g$ and its cohomology is the exterior algebra generated by the classes $\{\tau_j, \rho_j\}_{j=1}^g$.
Let
\[\omega=\sum_{j=1}^g \rho_j\tau_j\in H^2(\Pic^0(X)).\]
By rescaling the elements of the basis, we may assume that $\int_{\Pic^0(X)}\prod_{j=1}^g \rho_j\tau_j=\frac{1}{g!}$ so that $\omega^g\in H^{2g}(\Pic^0(X))$ is the class Poincaré dual to a point in $\Pic^0(X)$.

 Since $\ch_0(\BG_\Delta)=r$ and $\ch_1(\BG_\Delta)=q^\ast c_1(\Delta)$, we have
\[\xi_{\widetilde \BG}(\chh_1(\hat e_j))=r\tau_j \quad \textup{and}\quad \xi_{\widetilde \BG}(\chh_0(\hat f_j))=r\rho_j\]
in $H^\bullet\big(\widetilde M\big)=H^\bullet(M_\Delta\times \Pic^0(X))$ where we omit the obvious pullback. Let
\[W=\sum_{j=1}^g \chh_0(\hat f_j)\chh_1(\hat e_j)\in \BD^X\quad \textup{so that}\quad \xi_{\widetilde \BG}(W)=r^2\omega.\]
We will apply equation \eqref{eq: virasoromtilde} to descendents of the form $W^g D$ for some $D\in \BD^X$ and use that to deduce the Proposition. We have
\begin{align}\label{eq: virasoroMtilde2}
\source{virasoro-constraints-moduli-sheaves-vertex-algebras}
0&=\int_{[\widetilde M]^\vir}\xi_{\widetilde \BG}\big((\bL_k+\bS_k^\pt)(W^g D)\big)\\
&=r^{2g}\int_{[\widetilde M]^\vir}\omega^g \xi_{\widetilde \BG}\big((\bL_k+\bS_k^\pt)(D)\big)+\int_{[\widetilde M]^\vir}\xi_{\widetilde \BG}\big(\bR_k(W^g) D)\big)\nonumber\\
&=r^{2g}\int_{[M_\Delta]^\vir}\xi_{\BG_\Delta}\big((\bL_k+\bS_k^\pt)(D)\big)+gr^{2g-2}\int_{[\widetilde M]^\vir}\omega^{g-1}\xi_{\widetilde \BG}\big( \bR_k(W) D)\big).\nonumber
\end{align}
We now rewrite the second integral in terms of an integral in $M_\Delta$. We have
\[\xi_{\widetilde\BG}\big(\bR_k(W)\big)=(k+1)!\sum_{j=1}^g \xi_{\widetilde\BG}\big(\chh_0(\hat f_j)\chh_{k+1}(\hat e_j)\big)=r(k+1)!\sum_{j=1}^g \rho_j\xi_{\widetilde\BG}\big(\chh_{k+1}(\hat e_j)\big).\]
\begin{claim}
\source{virasoro-constraints-moduli-sheaves-vertex-algebras}
\notready
Let $h\colon \widetilde M\to M_\Delta$ be the projection and let $D\in \BD^X$. Then we have
\[h_\ast\big(g\omega^{g-1}\rho_j\xi_{\widetilde \BG}(D)\big)=\xi_{\BG_\Delta}\big(\bR_{-1}[e_j]D\big).\]
\end{claim}\begin{proof}[Proof of claim]
Let $J\subseteq H^\bullet\big(\widetilde M\big)$ be the annihilator ideal of $\omega^{g-1}\rho_j$; in particular, $H^{\geq 2}(\Pic^0(X))\subseteq J$. By definition,
\begin{align*}
\xi_{\widetilde \BG}\big(\chh_\bullet(\gamma)\big)&=p_\ast\big(\ch(\BG_\Delta)\ch(\CP)q^\ast \gamma\big)=p_\ast\big(\ch(\BG_\Delta)\big(1+(p^\ast \tau_j) (q^\ast e_j)+\ldots\big)q^\ast \gamma\big)\\
&=\xi_{\BG_\Delta}\big(\chh_\bullet(\gamma)\big)+\tau_j\xi_{\BG_\Delta}\big(\chh_\bullet(e_j\gamma)\big)+\ldots
\end{align*}
where the terms in $\ldots$ are all in the ideal $J$ and where we omit all the pullbacks via the projections $\widetilde M\to M_\Delta$ and $\widetilde M\to \Pic^0(X)$. Thus,
\begin{align*}
\xi_{\widetilde \BG}(D)=\xi_{\BG_\Delta}(D)+\tau_j\xi_{\BG_\Delta}\big(\bR_{-1}[e_i]D\big)+\ldots
\end{align*}
Since 
\[\int_{\Pic^0(X)}\omega^{g-1}\rho_j\tau_j=(g-1)!\int_X \prod_{l=1}^g \rho_l\tau_l=\frac{1}{g}\]
the claim is proven.\qedhere
\end{proof}
Given the claim, \eqref{eq: virasoroMtilde2} becomes
\[
0=r^{2g}\int_{M_\Delta}\xi_{\BG_\Delta}\big((L_k+S_k^\pt)(D)\big)+r^{2g-1}(k+1)!\int_{M_\Delta}\sum_{j=1}^g\xi_{\BG_\Delta}\big(R_{-1}[e_j]\chh_{k+1}(\hat e_j) D)\big)
\]
Since $h^{0, q}=0$ for $q>1$, we have
\[\bS_k=-r\bS_k^\pt-\sum_{j=1}^g \bR_{-1}[e_j]\chh_{k+1}(\hat e_j) \]
so the Proposition follows. \qedhere
\end{proof}
\begin{remark}
\source{virasoro-constraints-moduli-sheaves-vertex-algebras}
\notready\label{rem: fixed det}
\source{virasoro-constraints-moduli-sheaves-vertex-algebras}
When $S$ is a surface with $p_g>0$ we do not know a rigorous interpretation for the terms of the $\bS_k$ with $\gamma_t^L\in H^{0,2}(S)$. Heuristically, these should be related to a diagram like \eqref{eq: cartesianpic} where the Picard group is interpreted as a derived stack, see \cite{STV}. 
\end{remark}

\begin{example}
\source{virasoro-constraints-moduli-sheaves-vertex-algebras}
\notready\label{ex: rank2}
\source{virasoro-constraints-moduli-sheaves-vertex-algebras}
Let $M=M_C(2, \Delta)$ be the moduli space of stable bundles on a curve $C$ of genus $g$ with rank $2$ and fixed determinant $\Delta\in \Pic(C)$ of odd degree; this is a smooth moduli space of dimension $3g-3$. This moduli space has been studied a lot in the past and in particular its ring structure and descendent integrals are completely understood. We use this example to illustrate what kind of information the Virasoro constraints provide on descendent integrals. 
Let $\{e_1, \ldots, e_g\}\subseteq H^{0,1}(C)$ and $\{f_1, \ldots, f_g\}\subseteq H^{1,0}(C)$ be dual basis. Newstead proved in \cite[Theorem 1]{newstead} that the cohomology $H^\bullet(M)$ is generated by the classes
\[\eta=-2\xi\big(\chh_1(\1)\big), \quad \theta=4\xi\big(\chh_2(\pt)\big), \quad \xi_j=-\xi\big(\chh_1(e_j)\big), \quad \psi_j=\xi\big(\chh_2(f_j)\big).\footnote{In \cite{thaddeus, newstead} the classes $\eta, \theta, \xi_j, \psi_j, \zeta$ are called, respectively, $\alpha, \beta, \psi_{j+g}, \psi_{j}, \gamma$.}\]

The geometric realization $\xi$ is taken with respect to the sheaf $\BG\otimes \det(\BG)^{-1/2}$ (see Remark \ref{rem: normalizedpt}). Every descendent can be written explicitly in terms of the classes $\eta, \theta, \xi_j, \psi_j$. By \cite[(24) Proposition]{thaddeus}, one can then write every descendent integral in terms of integrals of products of the classes $\eta, \theta$ and
\[\zeta=2\sum_{j=1}^g\psi_j \xi_j.\]
These integrals are fully determined in \cite[(30)]{thaddeus}: for $m, k, p$ such that $m+2k+3p=3g-3$ we have
\begin{equation}\label{eq: integralsM21}
\source{virasoro-constraints-moduli-sheaves-vertex-algebras}
\int_M \eta^m \theta^k\zeta^p=(-1)^{g-1-p}\frac{m!g!}{q!(g-p)!}2^{2g-2-p}(2^q-2)B_q
\end{equation}
where $q=m+p-g+1$ and $B_q$ is a Bernoulli number. A careful combinatorial analysis shows that the Virasoro constraints given by Proposition \ref{prop: fixeddet} for $M$ are equivalent to the relations
\[(g-p)\int_M \eta^m \theta^k\zeta^p=-2m\int_M \eta^{m-1} \theta^{k-1}\zeta^{p+1}\]
which of course follows from \eqref{eq: integralsM21}; it would be interesting to have a direct and simpler proof of this identity. Note that the Virasoro constraints do not capture the most interesting part of \eqref{eq: integralsM21} which is the Bernoulli number. In some sense this is to be expected: since the Virasoro constraints hold in great generality (in particular are invariant under wall-crossing), it should not be expected that they can capture special information about particular moduli spaces. 
\end{example}
